\chapter{Conclusiones y trabajo futuro}
\label{cap:conclusiones}

\section{Conclusiones}
\label{sec:conc_principales}

Este Trabajo Fin de Máster ha desarrollado e integrado \framework, un
\emph{framework} bioinformático reproducible que combina modelos de
lenguaje y aprendizaje automático para la identificación de biomarcadores
transcriptómicos en cáncer de pulmón. Las conclusiones principales que se
extraen del trabajo son las siguientes:

\begin{enumerate}
  \item \textbf{Se ha demostrado la viabilidad de la curación clínica
  automatizada mediante LLM.} Llama~3.3-70b, servido vía Groq, alcanza
  una tasa de éxito de curación del 85,9\,\% de media sobre 11 cohortes
  GEO heterogéneas (figura~\ref{fig:curacion}), con nueve cohortes por
  encima del 95\,\%. Esto habilita la integración multi-cohorte a un
  coste marginal accesible y sin intervención manual.

  \item \textbf{La firma génica validada por consenso multi-cohorte es
  robusta.} Sobre la tarea de subtipo histológico ADC \emph{vs.} SQC, la
  firma alcanza AUC media LODO de 0,968 sobre tres cohortes
  independientes con balanced accuracy 0,940 (tabla~\ref{tab:lodo_sub}).
  Este rendimiento se mantiene en validación externa, sin depender de
  particiones internas del conjunto de entrenamiento.

  \item \textbf{Un panel mínimo clínicamente manejable de 20 genes
  aproxima el rendimiento de la firma completa.} Con AUC 0,966 frente al
  0,970 de la firma completa (1174 genes), la pérdida es inferior a
  0,01 (figura~\ref{fig:panel_minimo}). El panel es susceptible de ser
  medido mediante PCR cuantitativa o \emph{NanoString} y podría
  complementar la IHC diagnóstica en casos histológicamente ambiguos.

  \item \textbf{La firma recupera el panel diagnóstico IHC de la OMS
  con altísima cobertura.} De los 20 marcadores clínicos declarados en
  el anexo~\ref{anexo:paneles}, el framework recupera 18 (12/12
  escamosos, 6/8 adenocarcinoma) sin haberlos declarado durante el
  entrenamiento. Los dos marcadores no recuperados (SFTPA1, SFTPC) son
  proteínas del surfactante cuya expresión cae de forma homogénea entre
  ambos subtipos NSCLC.

  \item \textbf{La firma admite una interpretación biológica coherente.}
  Se identifican dos ejes reales: pérdida de arquitectura alvéolo-capilar
  normal (correlación media $\rho=-0{,}626$ con marcadores de alvéolo
  sano, hasta $\rho=-0{,}858$ en GSE31210, figura~\ref{fig:composicion})
  y actividad proliferativa aumentada. La firma no es una caja negra:
  refleja mecanismos oncológicos conocidos.

  \item \textbf{El clasificador tumor \emph{vs.} sano funciona bien tras
  recalibración.} Sobre 8 cohortes evaluables, AUC media LODO
  \textbf{0,925} y balanced accuracy 0,772 con umbral 0,5
  (tabla~\ref{tab:lodo_tvs}, figura~\ref{fig:auc_vs_balacc}). Aplicando
  recalibración isotónica anidada intra-train, la balanced accuracy sube a
  \textbf{0,810} (sección~\ref{sec:res_recalibracion}), con
  sensibilidad~0,988 y especificidad~0,632 ya equilibradas. La comparativa
  bajo el mismo LODO frente a Random Forest y SVM lineal
  (tabla~\ref{tab:comparativa_ml}) confirma que los tres modelos están en
  el mismo régimen; la elección de LASSO como modelo principal se
  justifica por interpretabilidad, no por rendimiento. Los IC 95\,\% por
  bootstrap (tabla~\ref{tab:ic_bootstrap}) sitúan la AUC
  \emph{pooled} en 0,943 [0,925; 0,959].

  \item \textbf{El panel se transfiere a RNA-Seq.} Aplicado a TCGA-LUAD
  (n=230) y TCGA-LUSC (n=178) —plataforma completamente distinta a la de
  entrenamiento— la firma alcanza AUC \textbf{0,857 [0,812; 0,899]}
  (sección~\ref{sec:res_tcga}, tabla~\ref{tab:tcga}). Ninguna muestra
  TCGA participó en la selección de genes ni en la calibración del
  modelo; la caída respecto al AUC LODO en microarray (0,968) es la
  esperable por el cambio tecnológico y no cuestiona la transferibilidad.

  \item \textbf{El framework es reproducible.} Todo el análisis se
  puede regenerar desde la raíz del repositorio con una única secuencia
  de comandos (anexo~\ref{anexo:despliegue}). Las semillas aleatorias
  están fijadas y los resultados intermedios se persisten en disco. La
  interfaz web (anexo~\ref{anexo:interfaz}) y el asistente
  conversacional (anexo~\ref{anexo:prompts}) leen los mismos ficheros,
  lo que garantiza la coherencia de las cifras mostradas.
\end{enumerate}

\section{Aportaciones originales}
\label{sec:conc_aportaciones}

Las aportaciones originales del trabajo, en orden de importancia, son:

\begin{itemize}
  \item Una \emph{tubería end-to-end reproducible} que integra por primera
  vez, hasta donde el autor tiene conocimiento, curación clínica basada
  en LLM, integración multi-cohorte, clasificación supervisada con
  validación externa LODO y verificación cruzada contra un panel clínico
  establecido.
  \item Una \emph{firma de 1174 genes} para el subtipo histológico de
  NSCLC validada estrictamente en tres cohortes independientes, y un
  \emph{panel mínimo de 20 genes} que aproxima su rendimiento.
  \item Un \emph{método de auditoría continua de coherencia} entre las
  cifras mostradas en la interfaz web (dashboard, entradilla, buscador de
  gen, asistente conversacional) y los datos subyacentes, que evita el
  problema clásico de cifras hardcodeadas que envejecen mal.
  \item Un \emph{asistente conversacional acotado a los resultados} con
  reglas explícitas para no inventar cifras ni citar biomarcadores que
  no formen parte de la firma, que ilustra un patrón de despliegue seguro
  de LLMs en herramientas científicas.
\end{itemize}

\section{Limitaciones}
\label{sec:conc_limitaciones}

Las principales limitaciones del trabajo son:

\begin{itemize}
  \item El entrenamiento se restringe a microarrays de Affymetrix. La
  transferencia directa a RNA-Seq no está evaluada.
  \item El clasificador de subtipo se limita a ADC y SQC. Para triage
  sobre casos sin diagnosticar es necesario ampliar el entrenamiento con
  otras histologías (neuroendocrinos, carcinoma de células grandes,
  metástasis de otros orígenes).
  \item El umbral del clasificador tumor \emph{vs.} sano requiere
  calibración por cohorte para uso clínico.
  \item Dos cohortes evaluables tienen tamaño muestral muy reducido
  ($n \leq 15$), lo que introduce ruido en las cifras individuales.
  \item El framework depende de servicios externos (API de Groq) que
  pueden cambiar sus tarifas o políticas en el futuro.
\end{itemize}

\section{Trabajo futuro}
\label{sec:conc_futuro}

Las líneas de trabajo futuro naturales incluyen:

\begin{description}
  \item[Validación prospectiva ampliada.] La transferibilidad a RNA-Seq
  ya se ha probado en este trabajo sobre TCGA-LUAD (n=230) y TCGA-LUSC
  (n=178) con AUC 0,857 [0,812; 0,899] (sección~\ref{sec:res_tcga}). El
  siguiente paso natural es una validación clínica real prospectiva:
  aplicar el panel a un lote de biopsias frescas procesadas por
  \emph{NanoString} o RT-qPCR y comparar con la histología clínica de
  cada caso.
  \item[Ampliación a subtipos neuroendocrinos.] Incorporar cohortes con
  tumores neuroendocrinos (LCNE, carcinoide, microcítico) y ampliar el
  clasificador de subtipo a una tarea multi-clase 4-5 categorías. Es el
  paso necesario para plantearse triage diagnóstico sobre casos sin
  diagnosticar.
  \item[Sustitución del LLM externo.] Explorar variantes locales
  (Llama~3.3-70b instalable, Phi-3, etc.) para eliminar la dependencia
  de la API externa. Es viable en máquinas con al menos 40\,GB de VRAM;
  para volúmenes de 1000-2000 muestras al mes puede ser más económico
  que la API.
  \item[Extensión metodológica a otros tumores.] Aplicar el mismo
  framework a patologías con datos GEO abundantes (mama, colon,
  hepatocarcinoma). El paso más costoso es adaptar los criterios de
  inclusión y los paneles IHC de validación a cada patología; el resto
  del pipeline es genérico.
  \item[Aplicación a firmas pronósticas.] Extender el framework para
  predecir supervivencia (regresión de riesgos proporcionales de Cox) en
  lugar de sólo clasificar histología. Requiere cohortes con datos de
  seguimiento clínico, disponibles en TCGA y en algunas cohortes GEO.
  \item[Publicación y despliegue.] Preparar un artículo científico y una
  versión pública de la interfaz que permita a otros investigadores
  aplicar el panel a sus propios datos.
\end{description}

\section{Balance personal}
\label{sec:conc_balance}

Desde el punto de vista formativo, el trabajo ha permitido al autor
integrar conocimientos de biología del cáncer, estadística, aprendizaje
automático, ingeniería de datos y desarrollo de aplicaciones web
interactivas en un único proyecto coherente. La colaboración con
modelos de lenguaje como asistentes de programación y como componentes
del propio pipeline ha sido una experiencia particularmente enriquecedora
y refleja, en pequeña escala, la transformación que estas herramientas
están introduciendo en la práctica bioinformática contemporánea.

El código del framework, la interfaz web y esta memoria se ponen a
disposición pública bajo licencia MIT, con la expectativa de que puedan
ser útiles a otros estudiantes de máster o investigadores que aborden
problemas similares.
